\chapter{Introducere}
\label{cap:introducere}

\section{Context și motivație}

Aproximativ unul din 36 de copii primește astăzi un diagnostic din spectrul autist, iar prevalența ADHD în populația pediatrică depășește 5\%. În spatele acestor cifre stau familii care caută intervenții timpurii, terapeuți supraîncărcați și — în România în mod special — un decalaj considerabil între nevoia de servicii specializate și disponibilitatea lor. Orice instrument care poate extinde capacitatea unui terapeut, fără a-i lua locul, are așadar o miză practică reală.

Două direcții de cercetare, dezvoltate independent, sugerează o astfel de oportunitate. Prima este terapia bazată pe LEGO (LEGO-Based Therapy), formalizată de dr. Daniel LeGoff \cite{legoff2004} pornind de la o observație simplă: copiii cu autism, care evită de regulă interacțiunile sociale, colaborează spontan atunci când mediul de joc este un set de piese LEGO — un sistem cu reguli fizice clare, predictibil și lipsit de ambiguitatea care face interacțiunea umană obositoare pentru ei. Studiile ulterioare \cite{legoff2006, owens2008} au confirmat că abilitățile sociale exersate în acest cadru se mențin și se transferă. A doua direcție este robotica asistivă: o literatură consistentă \cite{scassellati2012, robins2005, cabibihan2013} documentează un fenomen aparent paradoxal — mulți copii cu autism interacționează mai ușor cu roboții decât cu oamenii, fiindcă robotul este predictibil, răbdător și nu transmite semnale sociale contradictorii. Mai mult, studii precum cel al lui Kim și colaboratorii \cite{kim2013} arată că prezența robotului crește interacțiunea copilului \emph{cu adulții din încăpere}: robotul funcționează ca mediator social, nu ca substitut.

Legătura cu practica nu a rămas doar teoretică. În timpul dezvoltării proiectului am avut ocazia să lucrez alături de psihoterapeuții Marius și Andreea Lumpan, în cadrul primei inițiative de LEGO-Based Therapy din România dedicată copiilor neurodivergenți \cite{lumpanagerpres2025}. Experiența lor mi-a permis să observ direct cum interacționează copiii cu materialul LEGO în context terapeutic — colaborarea, toleranța la frustrare, ritmul sesiunii, rolul adultului care intervine înainte de escaladare. Aceste observații au influențat direct deciziile de design ale TriSense: instrucțiuni scurte, feedback fără penalizare, ritualuri predictibile și pauze de reglare.

TriSense s-a născut la intersecția acestor două direcții. L-am pornit ca proiect personal pentru competiția World Robot Olympiad (categoria Proiect), iar în cadrul acestei lucrări l-am dezvoltat într-un sistem complet — întreaga construcție, de la hardware la cod, fiind realizată individual. Întrebarea de la care am plecat a fost pragmatică: se poate construi, din componente accesibile oricărei școli sau cabinet — un set LEGO educațional, două plăci de dezvoltare de câteva zeci de lei și un laptop obișnuit — un robot capabil să susțină activități terapeutice reale, cu dialog vocal, jocuri structurate și feedback obiectiv pentru terapeut?

\section{Problema abordată}

Formulată tehnic, problema centrală a lucrării este \emph{proiectarea și implementarea unui sistem cyber-fizic care să orchestreze, în timp real și fiabil, componente diferite — un hub LEGO cu resurse minimale, un microcontroler cu Wi-Fi, o cameră independentă și servicii de inteligență artificială în cloud — într-o experiență de interacțiune coerentă, predictibilă și adaptată copiilor neurodivergenți.} Dificultatea nu stă în niciuna dintre componente luată separat, ci în integrarea lor: hub-ul LEGO nu are rețea, placa ESP32 are rețea dar resurse de calcul limitate, iar modelele de inteligență artificială rulează în cloud, cu latențe variabile — în timp ce utilizatorul final este un copil pentru care orice întrerupere sau comportament imprevizibil poate compromite sesiunea.

\section{Obiectivele și contribuțiile lucrării}

Obiectivul general este realizarea unui sistem robotic asistiv funcțional cap-coadă, validat pe hardware real. Din el decurg obiectivele specifice și, totodată, principalele contribuții — de natură implementativă, întrucât sistemul există, funcționează și a fost testat pe hardware real:

\begin{itemize}
    \item o \textbf{arhitectură distribuită pe trei niveluri} (hub LEGO cu Pybricks, placă LMS-ESP32 cu MicroPython, aplicație Python pe PC), cu protocoale adecvate fiecărui tip de trafic: MQTT pentru control, canal serial LPF2 pentru legătura hub--ESP32 și fluxuri TCP pentru audio;
    \item un \textbf{lanț vocal complet} — captură audio pe robot, transcriere automată, generarea răspunsului cu un model de limbaj și sinteză vocală redată pe difuzor — cu latență și calitate acceptabile pentru un dialog natural;
    \item un \textbf{canal bidirecțional} peste LPF2 care transmite și apăsările butoanelor fizice ale hub-ului, transformând robotul într-o interfață de intrare pentru copil;
    \item un \textbf{set de activități terapeutice} inspirate din literatura de specialitate: recunoașterea emoțiilor, imitarea tiparelor de mișcare, jocul de construcție „Build the Model” cu verificare vizuală, exercițiul de respirație, jocuri de funcții executive (Stop \& Go, estimarea timpului, categorii rapide), povestea co-construită și o rutină de dans;
    \item o \textbf{metodă de verificare a construcțiilor LEGO} bazată pe descrieri în limbaj natural și un model multimodal zero-shot, cu prag calibrat empiric, extensibilă fără date de antrenare;
    \item \textbf{colectarea automată a metricilor de sesiune} (activități, durate, încercări, reușite) într-un format direct utilizabil de terapeut.
\end{itemize}

\section{Structura lucrării}

Capitolul~\ref{cap:preliminarii} parcurge fundamentele teoretice și studiile de specialitate, inclusiv poziționarea față de orientările psihologice majore. Capitolul~\ref{cap:arhitectura} prezintă arhitectura sistemului și fluxurile de date. Capitolul~\ref{cap:implementare} coboară la nivelul implementării — firmware, aplicația de pe PC, jocurile și componenta de viziune — materialul detaliat fiind reluat în anexe. Capitolul~\ref{cap:evaluare} descrie testarea, rezultatele și considerațiile etice, iar capitolul~\ref{cap:concluzii} sintetizează concluziile și direcțiile de dezvoltare. Anexa~\ref{anexa:cod} prezintă structura proiectului software; catalogul detaliat al jocurilor și fragmentele de cod sunt în anexa~\ref{an:cap:implementare}.
